\documentclass{article}

% ICML 2025 style
\usepackage{icml2025}

\usepackage{amsmath}
\usepackage{amssymb}
\usepackage{booktabs}
\usepackage{graphicx}
\usepackage{subcaption}
\usepackage{xcolor}
\usepackage{hyperref}
\usepackage{microtype}

% Submission: use \icmlfinalcopy to show author names
% \icmlfinalcopy

\begin{document}

\twocolumn[
\icmltitle{Distilling Depth: Improving a Compact Monocular Depth Model\\
           via Teacher Supervision and Novel View Synthesis}

\icmlsetsymbol{equal}{*}

\begin{icmlauthorlist}
\icmlauthor{Christian Deubel}{eth}
\icmlauthor{Noah Meurer}{eth}
% TODO: add remaining team members
\end{icmlauthorlist}

\icmlaffiliation{eth}{Department of Computer Science, ETH Z\"{u}rich, Z\"{u}rich, Switzerland}

\icmlcorrespondingauthor{Christian Deubel}{cdeubel@student.ethz.ch}

\icmlkeywords{monocular depth estimation, knowledge distillation, pseudo-labels, novel view synthesis, Gaussian splatting}

\vskip 0.3in
]

\printAffiliationsAndNotice{}

%% -----------------------------------------------------------------------
\begin{abstract}
% TODO: fill in after results are known
Large foundation models for monocular depth estimation achieve strong zero-shot performance but are prohibitively expensive for edge deployment.
We study whether a compact student model can be systematically improved by leveraging a larger teacher without requiring additional ground-truth annotations.
Starting from a zero-shot evaluation of the student (\emph{Baseline 1}), we explore three training strategies of increasing sophistication:
(i) fine-tuning on ground-truth depth labels with a frozen backbone (\emph{Baseline 2}),
(ii) replacing noisy ground-truth labels with pseudo-labels from a stronger teacher model (\emph{Baseline 3}), and
(iii) augmenting the training set with novel synthetic views rendered via teacher-guided Gaussian Splatting (\emph{Baseline 4}).
Experiments on the CIL monocular depth benchmark show that \ldots\ % TODO
\end{abstract}


%% -----------------------------------------------------------------------
\section{Introduction}
\label{sec:intro}

% TODO: cite MiDaS, DPT, ZoeDepth, Depth Anything v1/v2/v3, UniDepth, Marigold as related prior work
Monocular depth estimation — recovering per-pixel depth from a single RGB image — is a core problem in computer vision with applications ranging from autonomous driving to robotics and augmented reality~\cite{TODO_midas,TODO_dpt,TODO_depthany_v2}.
Recent large-scale foundation models have pushed the frontier considerably \cite{TODO_depthany_v3,TODO_unidepth,TODO_marigold}, but their parameter counts and compute requirements make them unsuitable for deployment on resource-constrained hardware.

A natural question is: \emph{can knowledge from a large teacher model be distilled into a compact student, closing the performance gap without the inference cost?}
Knowledge distillation \cite{TODO_hinton_kd} is a well-studied paradigm in classification, but its application to dense prediction tasks such as monocular depth is less explored.
Pseudo-labelling \cite{TODO_pseudolabel} offers a simple but effective avenue: the teacher annotates unlabelled (or noisily labelled) data, and the student is supervised on these richer targets.

Beyond pseudo-labelling, recent advances in 3D scene representations — in particular 3D Gaussian Splatting (3DGS) \cite{TODO_3dgs} — open a new dimension: given a single image and a depth estimate, it is possible to render plausible novel views with consistent depth supervision.
This is analogous to classical image augmentation (crops, flips, colour jitter) for discriminative models, but operates in 3D and can expose the student to viewpoint variation it would never see from a fixed single-view dataset.

In this paper we present a systematic study of these ideas through four baselines on the CIL 2026 monocular depth estimation benchmark:
\begin{enumerate}
    \item \textbf{Zero-shot student} — evaluate DA3MONO-LARGE out-of-the-box.
    \item \textbf{GT fine-tuning} — freeze the backbone and fine-tune the prediction head on provided ground-truth depth.
    \item \textbf{Pseudo-label distillation} — replace GT with teacher pseudo-labels from DA3NESTED-GIANT.
    \item \textbf{NVS augmentation} — expand the training set with novel synthetic views rendered by 3DGS and supervised by the teacher.
\end{enumerate}
Our goal is not to beat the teacher at inference time, but to understand how much of its knowledge can be transferred to a model that is practical to deploy.


%% -----------------------------------------------------------------------
\section{Related Work}
\label{sec:related}

% TODO: flesh out with proper citations after reference research

\paragraph{Monocular depth estimation.}
Early learning-based methods framed depth estimation as direct regression \cite{TODO_eigen2014}.
DPT \cite{TODO_dpt} introduced vision-transformer backbones with multi-scale feature fusion, substantially improving detail preservation.
MiDaS \cite{TODO_midas} demonstrated that training on diverse mixed datasets yields strong zero-shot generalisability.
The Depth Anything series \cite{TODO_depthany_v1,TODO_depthany_v2,TODO_depthany_v3} scaled this recipe further, combining large unlabelled image collections with metric depth supervision.
ZoeDepth \cite{TODO_zoedepth} and UniDepth \cite{TODO_unidepth} additionally recover metric-scale depth.
Marigold \cite{TODO_marigold} repurposed latent diffusion models for depth, achieving high-fidelity predictions at the cost of slow inference.

\paragraph{Knowledge distillation.}
Hinton et al.\ \cite{TODO_hinton_kd} introduced soft-target distillation for classification.
Feature-level distillation \cite{TODO_fitnets,TODO_relationkd} transfers intermediate representations.
For dense prediction, \cite{TODO_dense_kd} showed that pixel-wise distillation can compress segmentation models effectively.
Pseudo-labelling \cite{TODO_pseudolabel,TODO_noisy_student} is a related semi-supervised technique in which a teacher labels unlabelled data; the student then trains on both labelled and pseudo-labelled samples.

\paragraph{Novel view synthesis and 3D Gaussian Splatting.}
NeRF \cite{TODO_nerf} established volumetric scene representations for novel view synthesis.
3D Gaussian Splatting \cite{TODO_3dgs} achieves real-time rendering by representing scenes as learnable 3D Gaussians.
Several works have explored initialising 3DGS from monocular depth \cite{TODO_depth_gs,TODO_monogs}.
Using NVS as a data augmentation strategy for downstream recognition tasks has been explored in \cite{TODO_nvs_aug}, motivating our Baseline 4.

\paragraph{LoRA and parameter-efficient fine-tuning.}
LoRA \cite{TODO_lora} injects low-rank adapter matrices into frozen transformer weights, enabling efficient fine-tuning with a fraction of the parameters.
We apply LoRA to the DPT feature-extraction blocks of the student backbone as an alternative to full-head-only fine-tuning.


%% -----------------------------------------------------------------------
\section{Methodology}
\label{sec:method}

\subsection{Problem Setup}

Let $\mathcal{D}_{\text{train}} = \{(x_i, d_i)\}_{i=1}^{N}$ be the training set of RGB images $x_i \in \mathbb{R}^{H \times W \times 3}$ and corresponding metric depth maps $d_i \in \mathbb{R}^{H \times W}_{>0}$.
A held-out test set $\mathcal{D}_{\text{test}}$ provides unlabelled images for final evaluation.
All images are resized to $560 \times 560$ pixels.
Depth values are clipped to the range $[d_{\min}, d_{\max}] = [0.001\text{m},\, 80\text{m}]$ and normalised to $[0, 1]$ for training:
\begin{equation}
    \tilde{d} = \frac{\operatorname{clip}(d,\, d_{\min},\, d_{\max}) - d_{\min}}{d_{\max} - d_{\min}}.
\end{equation}

\paragraph{Student model.}
We use DA3MONO-LARGE, the large single-view variant of Depth Anything 3~\cite{TODO_depthany_v3}.
The architecture follows a DPT decoder \cite{TODO_dpt} on top of a Vision Transformer backbone \cite{TODO_vit}.
Features are extracted from backbone layers 4, 11, 17, and 23 and fused through four DPT blocks into a final head that produces a dense depth map.

\paragraph{Teacher model.}
Baselines 3 and 4 employ DA3NESTED-GIANT-LARGE-1.1, a larger nested model that jointly processes sequences of frames and achieves higher accuracy at the cost of greater compute.
For an upper-bound reference we also evaluate DA3-GIANT-1.1 zero-shot.

\paragraph{Training loss.}
All fine-tuning runs minimise the scale-invariant logarithmic (SiLog) loss \cite{TODO_eigen2014}:
\begin{equation}
    \mathcal{L}_{\text{SiLog}}(\hat{d}, d) = \sqrt{
        \frac{1}{|\mathcal{V}|}\sum_{p \in \mathcal{V}} \delta_p^2
        - \frac{\lambda}{|\mathcal{V}|^2}\!\left(\sum_{p \in \mathcal{V}} \delta_p\right)^{\!2}
    },
    \label{eq:silog}
\end{equation}
where $\delta_p = \log \hat{d}_p - \log d_p$, $\mathcal{V}$ is the set of valid pixels ($\hat{d}_p > 0$ and $d_p > 0$), and $\lambda = 0.5$.
SiLog is invariant to global scale shifts, making it robust to the scale ambiguity present in monocular depth.

\subsection{Baseline 1: Zero-Shot Student Evaluation}

We load the pre-trained DA3MONO-LARGE weights and run inference on $\mathcal{D}_{\text{test}}$ without any task-specific training.
Images are processed at $560 \times 560$ resolution using an upper-bound resize strategy (no padding artefacts).
This baseline quantifies the student's generalisation from its large-scale pre-training and sets the lower bound for our improvement study.

\subsection{Baseline 2: Fine-Tuning on Ground-Truth Labels}

Na\"ively fine-tuning all parameters on the small training set risks catastrophic forgetting of the rich depth priors acquired during large-scale pre-training.
We therefore investigate two parameter-efficient strategies.

\paragraph{Full-head fine-tuning (B2-Head).}
The ViT backbone is fully frozen; only the DPT head parameters are updated.
This preserves all pre-trained feature representations while adapting the decoder to the dataset's depth distribution.

\paragraph{LoRA on DPT blocks (B2-LoRA).}
As an alternative, we apply Low-Rank Adaptation (LoRA) \cite{TODO_lora} to the attention and feed-forward projections of the four DPT feature-extraction blocks (backbone layers 4, 11, 17, 23).
Concretely, we inject rank-$r = 4$ adapter matrices with scaling $\alpha = 16$ into the query/key/value and output projections (\texttt{qkv}, \texttt{proj}) and both MLP layers (\texttt{fc1}, \texttt{fc2}) of each targeted block.
The backbone operates in \texttt{eval()} mode throughout training (no dropout); only the LoRA parameters and the head are updated.

\paragraph{Training details.}
We train for $E = 3$ epochs with batch size 16, AdamW optimiser ($\text{lr} = 10^{-6}$, weight decay $10^{-2}$), cosine annealing schedule (min lr $= 0.01 \cdot \text{lr}$), and gradient clipping at norm $1.0$.
Mixed precision (bfloat16) is enabled via PyTorch AMP.
A random 10\% split of the training set is held out as a validation set using a fixed seed; the best checkpoint by validation SiLog is used for test inference.

\subsection{Baseline 3: Pseudo-Label Distillation}

Instead of supervising the student on the (potentially noisy) ground-truth depth, we replace the training targets with pseudo-labels generated by the stronger teacher model.

\paragraph{Pseudo-label generation.}
We run DA3NESTED-GIANT-LARGE-1.1 in zero-shot mode over all training images $\{x_i\}$, saving float32 depth maps at $560 \times 560$ resolution.
Invalid pixels (non-finite values or $d \leq 0$) are set to zero and excluded from the loss via the valid-pixel mask $\mathcal{V}$.

\paragraph{Student fine-tuning.}
The student is then fine-tuned on the pseudo-labelled dataset $\{(x_i, \hat{d}^T_i)\}$ using exactly the training protocol of Baseline 2 (head-only or LoRA).
The key difference from Baseline 2 is the supervision signal: teacher pseudo-labels provide richer spatial detail and are free of sensor noise, but may introduce their own systematic biases.

\subsection{Baseline 4: Novel View Synthesis Augmentation (In Progress)}

A single-view depth model never observes the same scene from a different viewpoint during training, limiting the geometric consistency it can learn.
Baseline 4 addresses this by augmenting the training set with synthetic novel views generated via 3D Gaussian Splatting.

\paragraph{Pipeline overview.}
\begin{enumerate}
    \item \textbf{Depth initialisation.} For each training image $x_i$, the teacher model produces a metric depth map $\hat{d}^T_i$ which is used to initialise a point cloud.
    \item \textbf{3DGS fitting.} A 3DGS scene is fit to the (image, depth) pair, yielding a compact set of 3D Gaussians that can be rendered from arbitrary viewpoints.
    \item \textbf{Novel view rendering.} Small random camera perturbations (translations up to $\delta_t$, rotations up to $\delta_\theta$) are applied to render augmented RGB images $x_i^{(k)}$ together with their corresponding rendered depth maps $d_i^{(k)}$.
    \item \textbf{Student training.} The student is fine-tuned on the union of original and augmented samples.
\end{enumerate}

The augmented depth targets are derived from the 3DGS representation (ground truth for the rendered geometry) rather than from a second teacher inference pass, making the labels internally consistent across viewpoints.
This scheme is analogous to spatial augmentations (crops, flips) for classification, but extends to geometrically coherent 3D transformations.

\paragraph{Expected benefit.}
By training on multiple synthetic viewpoints of each scene, the student observes more diverse spatial configurations and is encouraged to learn representations that are robust to small viewpoint changes — a useful prior for real-world deployment where camera pose is not perfectly controlled.


%% -----------------------------------------------------------------------
\section{Experiments}
\label{sec:experiments}

\subsection{Dataset}

All experiments use the CIL 2026 Monocular Depth Estimation benchmark.
The training set contains $N = \text{TODO}$ RGB–depth pairs; the test set contains $\text{TODO}$ images without ground-truth depth.
Depth maps are stored as float32 \texttt{.npy} arrays at $560 \times 560$ resolution.
% TODO: describe scene types (indoor/outdoor), sensor, scale range

\subsection{Evaluation Metrics}

Following standard practice in monocular depth estimation \cite{TODO_eigen2014}, we report the following metrics on the test set:
\begin{itemize}
    \item \textbf{SiLog} — scale-invariant log loss (Eq.~\ref{eq:silog}); our primary metric.
    \item \textbf{AbsRel} — mean absolute relative error: $\frac{1}{|\mathcal{V}|}\sum_{p} |d_p - \hat{d}_p| / d_p$.
    \item \textbf{$\delta_1$} — percentage of predictions with $\max(d_p/\hat{d}_p,\, \hat{d}_p/d_p) < 1.25$.
\end{itemize}

\subsection{Quantitative Results}

\begin{table}[h]
\centering
\caption{Quantitative comparison of all baselines on the CIL test set. Lower is better for SiLog and AbsRel; higher is better for $\delta_1$. B4 results pending implementation.}
\label{tab:main}
\begin{tabular}{lrrr}
\toprule
Method & SiLog $\downarrow$ & AbsRel $\downarrow$ & $\delta_1$ $\uparrow$ \\
\midrule
Teacher (DA3-GIANT, zero-shot)       & TODO & TODO & TODO \\
\midrule
B1: Student zero-shot                & TODO & TODO & TODO \\
B2-Head: GT fine-tune (head only)    & TODO & TODO & TODO \\
B2-LoRA: GT fine-tune (LoRA)         & TODO & TODO & TODO \\
B3: Pseudo-label distillation        & TODO & TODO & TODO \\
B4: NVS augmentation                 & TODO & TODO & TODO \\
\bottomrule
\end{tabular}
\end{table}

\subsection{Qualitative Results}

Figure~\ref{fig:qualitative} shows representative depth predictions from each baseline alongside the ground-truth depth and the input RGB image.

\begin{figure*}[t]
    \centering
    \begin{subfigure}[b]{0.16\linewidth}
        \includegraphics[width=\linewidth]{figures/qual_rgb.png}
        \caption{RGB input}
    \end{subfigure}
    \begin{subfigure}[b]{0.16\linewidth}
        \includegraphics[width=\linewidth]{figures/qual_gt.png}
        \caption{Ground truth}
    \end{subfigure}
    \begin{subfigure}[b]{0.16\linewidth}
        \includegraphics[width=\linewidth]{figures/qual_b1.png}
        \caption{B1 (zero-shot)}
    \end{subfigure}
    \begin{subfigure}[b]{0.16\linewidth}
        \includegraphics[width=\linewidth]{figures/qual_b2.png}
        \caption{B2 (GT fine-tune)}
    \end{subfigure}
    \begin{subfigure}[b]{0.16\linewidth}
        \includegraphics[width=\linewidth]{figures/qual_b3.png}
        \caption{B3 (pseudo-label)}
    \end{subfigure}
    \begin{subfigure}[b]{0.16\linewidth}
        \includegraphics[width=\linewidth]{figures/qual_teacher.png}
        \caption{Teacher (ref.)}
    \end{subfigure}
    \caption{Qualitative depth predictions. Colour mapping: blue = near, yellow = far (viridis). % TODO: select representative examples showing failure modes and improvements}
    \label{fig:qualitative}
\end{figure*}

\subsection{Ablation Study}

\begin{table}[h]
\centering
\caption{Ablation: effect of fine-tuning strategy in Baseline 2.}
\label{tab:ablation_b2}
\begin{tabular}{lrr}
\toprule
Fine-tune scope & SiLog $\downarrow$ & AbsRel $\downarrow$ \\
\midrule
None (B1, zero-shot)  & TODO & TODO \\
Head only (B2-Head)   & TODO & TODO \\
LoRA $r$=4 (B2-LoRA)  & TODO & TODO \\
\bottomrule
\end{tabular}
\end{table}


%% -----------------------------------------------------------------------
\section{Discussion}
\label{sec:discussion}

% TODO: fill in after results; prompts below

% Does fine-tuning on GT improve over zero-shot? If not, why (dataset size, label noise)?
% Does pseudo-label supervision from the teacher outperform GT supervision?
% What is the gap between teacher zero-shot and best student after distillation?
% Does NVS augmentation provide additional gains over pseudo-label-only training?
% Where does the student still fail qualitatively (thin structures, reflective surfaces, sky)?


%% -----------------------------------------------------------------------
\section{Conclusion}
\label{sec:conclusion}

% TODO: fill in after results

We studied four progressive strategies for improving a compact monocular depth student model using supervision from a larger teacher.
Our results show that \ldots\ % TODO
In particular, pseudo-label distillation \ldots\ % TODO
Novel view synthesis augmentation \ldots\ % TODO
These findings suggest that teacher-guided training is a practical and label-efficient path to deploying accurate depth estimation on resource-constrained hardware.


%% -----------------------------------------------------------------------
\bibliography{references}
\bibliographystyle{icml2025}

\end{document}
